\chapter{凸領域}
    \section{凸領域内のベクトルの凸結合は元の凸領域に属す}
        \begin{shadebox}
            $V$を凸領域とし，$\vv_1,\dotsc,\vv_n\in V,\alpha_1,\dotsc,\alpha_n\geq 0,\sum_{i=1}^n\alpha_i = 1$であるとき，$\sum_{i=1}^n\alpha_i\vv_i \in V$
        \end{shadebox}
        \begin{proof}
            $n=1$のときは明らか。
            $n=m\in\naturalNumbers$まで成り立つと仮定する。
            $n=m+1$のとき，
            \[ \sum_{i=1}^n\alpha_i\vv_i = \sum_{i=1}^{m+1}\alpha_i\vv_i = \sum_{i=1}^m\alpha_i\vv_i + \alpha_{m+1}\vv_{m+1} = A_m\sum_{i=1}^m\beta_i\vv_i + \alpha_{m+1}\vv_{m+1} \quad \parens*{A_m \coloneq \sum_{i=1}^m\alpha_i,\beta_i \coloneq \frac{\alpha_i}{A_m}} \]
            $\sum_{i=1}^m\beta_i = 1$であるから，仮定より$\bm{w}_{m+1} \coloneq \sum_{i=1}^m\beta_i\vv_i\in V$である。
            そして$A_m\geq 0,A_m + \alpha_{m+1} = 1$であるから$\sum_{i=1}^n\alpha_i\vv_i = A_m\bm{w}_{m+1} + \alpha_{m+1}\vv_{m+1} \in V$である。
        \end{proof}
    \section{相対錐の共通部分の相対錐}
        \begin{shadebox}
            $\vv_1,\dotsc,\vv_m\in\realNumbers^n \;(m\leq n)$は一次独立であるする。
            $\vv\in\realNumbers^n$の相対錐は次式で定義される。
            \[P(\vv)\coloneq\braces*{\bm{w}|\inProd{\vv}{\bm{w}}\geq0}\]
            で定義する。あるベクトル$\bm{u}$が$\bigcap_{i=1}^{m}{P(\vv_i)}$の相対錐にある，すなわち
            \[\inProd{\bm{u}}{\bm{w}}\geq0,\forall\bm{w}\in\bigcap_{i=1}^{m}{P(\vv_i)}\]
            であるとき，$\bm{u}$は適当な定数$a_i\geq0(i=1,\dotsc,m)$によって
            \[\bm{u} = \sum_{i=1}^{m}{a_i\vv_i}\]
            と表せる。
        \end{shadebox}
        \begin{proof}
            まず$\bm{u}$が$\vv_1,\dotsc,\vv_m$の一次結合で表せることを示す。\\
            $W \coloneq \vecSpan\bracks*{\vv_1,\dotsc,\vv_m}$とすると$\realNumbers^n = W\oplus W^\bot$であるから，適当な定数$c_1,\dotsc,c_m$及び適当なベクトル$\bm{w}_t\in W^\bot$を用いて
            \begin{equation}
                \bm{u} = \sum_{i=1}^{m}{c_i\vv_i} + \bm{w}_t
            \end{equation}
            と表せる。$\bm{p}$を$\bigcap_{i=1}^{m}{P(\vv_i)}$の任意のベクトルとすると
            \[\inProd{\vv_i}{\bm{p}}\geq0,\forall\vv_i\in W\]
            また，$\bm{w}_{t2}$を$W^\bot$の任意のベクトルとすると
            \[\inProd{\vv_i}{\bm{w}_{t2}}=0,\forall\vv_i\in W\]
            従って$\bm{p}' \coloneq \bm{p}+\bm{w}_{t2}$とすると
            \[\inProd{\vv_i}{\bm{p}'}\geq0,\forall\vv_i\in W\]
            であるから$\bm{p'}$もまた$\bigcap_{i=1}^{m}{P(\vv_i)}$のベクトルである。従って$\inProd{\bm{u}}{\bm{p}'}\geq0$であるから
            \[0\leq\inProd{\bm{u}}{\bm{p}'} = \inProd{\sum_{i=1}^{m}{c_i\vv_i}}{\bm{p}} + \inProd{\bm{w}_{t1}}{\bm{w}_{t2}}\]
            である。然るに$\bm{w}_{t2}\in W^\bot$は任意であるから，$\bm{w}_{t1}=\bm{0}$でない限り$\inProd{\bm{w}_{t1}}{\bm{w}_{t2}}$の値をいくらでも小さくできる。
            よって上式が成り立つためには$\bm{w}_{t1}=\bm{0}$でなくてはならならず，式(1)より結局
            \[\bm{u} = \sum_{i=1}^{m}{c_i\vv_i}\]
            と表せることになる。
            \par
            次に$c_i$が全て非負であることを示す。%そのために$\parens*{\bm{d}_i\in\bigcap_{i=1}^{m}{P(\vv_i)}}\land\inProd{\vv_j}{\bm{d}_i}=\delta_{ij}\;(\delta_{ij}\text{はKroneckerのデルタ})$なるベクトル$\bm{d}_i$が存在することを示そう。
            $\vv_1,\dotsc,\vv_m\;(m\leq n)$は一次独立であったから次の行列
            \[
                A = \bracks*{
                    \begin{array}{c}
                        {\vv_1}^\mathrm{T}\\
                        {\vv_2}^\mathrm{T}\\
                        \vdots\\
                        {\vv_m}^\mathrm{T}
                    \end{array}
                }
            \]
            の階数は$m$である。従って次の方程式
            \[
                A\vx = \bracks*{
                    \begin{array}{c}
                        0\\
                        \vdots\\
                        1\\
                        \vdots\\
                        0
                    \end{array}
                }
                \begin{array}{c}
                    \quad\\
                    \quad\\
                    \leftarrow i\\
                    \quad\\
                    \quad
                \end{array}
            \]
            は解を持つ。この解を$\bm{d}_i$とすれば$\inProd{\vv_j}{\bm{d}_i}=\delta_{ij}$であるから$\bm{d}_i\in\bigcap_{i=1}^{m}{P(\vv_i)}$。
            よって
            \[0\leq\inProd{\bm{u}}{\bm{d}_i} = c_i\]
        \end{proof}